\subsection{Symbolic Program Construction and Correspondence}
\label{app:symbolic-program-construction-and-correspondence}
We prove Theorem~\ref{thm:symbolic-compilation}: the generated program gives
Section~\ref{sec:attribute-synthesis-algorithm-and-guarantees}'s feasibility result and, when
feasible, the same RHS attributes. We establish this correspondence in
four steps. Correctness of the mathematical construction itself is already
proved in Section~\ref{sec:correctness-and-completeness}.

\begin{proof}[Proof of Theorem~\ref{thm:symbolic-compilation}]
Fix an admitted plan satisfying Section~\ref{sec:formal-scope-and-guarantees}, and let $m$
be arbitrary legal match metadata in that domain. The frontend has already
identified the operation templates, captured source fields, and independent
RHS axis unknowns. As in Section~\ref{sec:symbolic-attribute-builder-construction},
$E_{\mathrm{check}}$ is the compiled feasibility check and $E_{\mathrm{build}}$ is the compiled builder.

\paperparagraph{Step 1: Preserve logical axes and their connections.}
The generated expressions must first describe the same axes as the
mathematical model. Each input dimension starts with its own identity.
Following the LHS operations, a permutation reorders axes, a sum removes
axes from the output, and a contraction connects its declared batch and
contracting partners. Removing an axis from the output preserves its
recorded connections. Each operation also records its declared output
order. Induction over the LHS therefore gives the same logical-axis
classes and ordered input and output lists as Section~\ref{sec:axis-based-representation}.
Equal extents do not create connections, and additive components retain
separate axis identities.

Mappings and partitions preserve these identities. Gather composition
selects the same entry $Q[p[h[j]]]$ on both sides of
Equation~\eqref{eq:gather-composition}. The partition selectors satisfy
$\gath(c_L,u_G)=\gath(p_G,v_G)$. Applying $u_G$ to $c_R$ preserves
their partners. The selected positions and their ordered complement cover
each group without overlap. Both selection order and remaining-group order
are recorded explicitly. Thus subsequent passes use the same axes and
pairings as the mathematical construction, for every finite group length.

\paperparagraph{Step 2: Produce the same feasibility result.}
The next task is to determine whether the compiled feasibility check makes exactly
the same decision as $\operatorname{check}(m)$. The inputs $\mathrm{In}_v$ inside each
RHS subtree are known at compile time. The generated expressions identify
an axis as mandatory when it occurs inside that subtree and is needed in the final
output or by an outside input:
\begin{equation}
 M_v=\left(\bigcup_{X\in \mathrm{In}_v}\operatorname{set}(I_X)\right)
 \cap\left(\operatorname{set}(\operatorname{Ord}_{\mathrm{out}})\cup
             \bigcup_{X\notin \mathrm{In}_v}\operatorname{set}(I_X)\right).
 \label{eq:emitted-mandatory}
\end{equation}
This is the mandatory-axis condition of Equation~\eqref{eq:mandatory-axes}.
The forward pass then applies the same retained-axis rules: inputs keep
their axes, permutations keep the child's set, sums intersect it with
$M_v$, and contractions use Equation~\eqref{eq:canonical-roles}. Induction
from inputs to root therefore gives the same $S_v$ at every node.

Starting with $\operatorname{Ord}_{\mathrm{out}}$, the reverse pass computes the same requirements $Req_v$
and local ordering checks as Section~\ref{sec:checking-axis-order}.
The checks require distinct available axes, the input subsequence condition,
and the declared batch/left-free/right-free order at contractions.
The root check requires $S_{\mathrm{root}}=\operatorname{set}(\operatorname{Ord}_{\mathrm{out}})$.
Set enumeration cannot affect these results: membership tests ignore order,
while ordered filters preserve their input order.

These checks use only finite set and list operations, with no partial
position lookup. They are evaluated before the attribute computation, so
a failed check does not require evaluating a builder. Their conjunction
is exactly Equation~\eqref{eq:feasibility}. For multiple terms, the component
checks are combined as in Appendix~\ref{app:extension-to-multiple-terms}. Hence $E_{\mathrm{check}}$
returns the same Boolean result as $\operatorname{check}(m)$, including
when construction is infeasible. This proves Equation~\eqref{eq:emitted-check}.

\paperparagraph{Step 3: Produce the same attributes when the check succeeds.}
Assume the check succeeds. We must now show that the generated position
queries are defined and choose the same attributes as the mathematical
builder. Induct over RHS nodes using Appendix~\ref{app:operation-construction-rules}.
Inputs use their original axis orders. A permutation places $Req_v$ first
and appends the remaining axes in child order. A sum removes the selected
axes and preserves the survivors' order.

At a contraction, the batch order starts with the required batch axes,
then appends the remaining batch axes in left-operand order. Contracting
axes also follow left-operand order. Both operands use the same logical-axis
sequences. These choices are exactly the ordering rule $\tau$ used by
$\operatorname{build}_{\tau}$.

By induction, each child order lists its retained axes exactly once.
The role sets and successful checks ensure that every selected axis is
present in its operand. Every position query is therefore defined;
source legality supplies equality of paired extents. Since the generated
program uses the same axis sets, output orders, and position queries as the
mathematical construction, it produces the same attributes. Each node's
output contains exactly $S_v$ and satisfies $Req_v$, and the root has order
$\operatorname{Ord}_{\mathrm{out}}$. Independent component assignments combine as in
Appendix~\ref{app:extension-to-multiple-terms}. This proves Equation~\eqref{eq:emitted-build}.

\paperparagraph{Step 4: Construct a finite program without concrete ranks.}
Finally, the compiler must construct this program before match metadata
is available. The admitted plan and its component expansion are finite,
and each pass appends finitely many expressions per node. Each expression
refers only to previously constructed expressions, which may be reused.
Variable-length axis lists are operands of these expressions; compilation
does not enumerate their concrete entries. The expressions are chosen from
the plan and operation rules without consulting the written guard. Thus
the compiler constructs a finite program without concrete ranks or
successful symbolic simplification. This completes the correspondence
and termination arguments.
\end{proof}
